%% kapa: exam=1 round=19 year=2008 period=1 form=A subject=appraisal_law subject_ko=부동산관계법규 seq=3 number=43
%% source: Q-Net 감정평가사 기출문제 (cst00309) artlSeq=1001962 file=extracted/1008/1001962/2008년 제19회 감정평가사 1차시험 기출문제/1교시.hwp
%% conversion: kordoc 4.24.51 -> tools/gen1.py (auto); review=proofread-sonnet; flags=none
\begin{question}{43}
국토의 계획 및 이용에 관한 법령상 개발제한구역에 관한 설명으로 옳지 않은 것은?
\begin{choices}
  \item 국방부장관의 요청에 의해 보안상 도시의 개발을 제한할 필요가 있다고 인정되는 경우에도 개발제한구역이 지정될 수 있다.
  \item 개발제한구역의 지정 또는 변경에 관한 도시관리계획은 국토해양부장관이 결정한다.
  \item 개발제한구역에서 해제되는 구역으로 계획적 개발 또는 관리가 필요한 지역의 전부 또는 일부에 대하여 제1종지구단위계획구역을 지정할 수 있다.
  \item 시․도지사는 개발제한구역이 해제되는 지역에 대하여 해제 이후 최초로 결정되는 도시관리계획을 결정함에 있어서 미리 국토해양부장관과 협의하여야 한다.
  \item 개발제한구역 안의 취락을 정비하기 위하여 지정되는 용도지구는 자연취락지구이다.
\end{choices}
\end{question}
\answer{5}
